\chapter{总结与展望}\label{chap:conclusion}

\section{工作总结}

本文围绕动态非合作目标抓取问题，实现并评估了一套基于合成专家数据的扩散模仿学习控制流程。系统首先在仿真中构建中等难度动态立方体任务，再利用有限状态机和MPC专家生成示教数据，随后将成功回合编译为目标中心化训练窗口，训练条件扩散策略输出末端执行器任务空间参考。运行时，扩散策略负责抓取前阶段的动作块生成，MPC控制器负责跟踪、约束处理和后续运输。

实验结果表明，在真值观测完整中等场景中，扩散策略达到85.0\%成功率，低于MPC专家的95.0\%；但在若干分解运动场景中，扩散策略已接近或达到专家水平。含噪实验进一步表明，扩散策略在中等高斯噪声和强噪声组合下显著优于直接使用噪声观测的MPC专家，说明历史观测、动作块生成和时间集成对高频抖动具有一定抑制作用。与此同时，扩散策略在轻噪声、漂移、离群点和丢帧延迟条件下仍落后于MPC专家，说明当前模型尚未形成全面鲁棒性。

\section{主要限制}

本文仍存在以下限制。首先，全部实验均在仿真中进行，尚未验证真实传感器、真实机械臂和真实接触动力学下的表现。其次，策略训练主要基于真值专家数据，噪声分布覆盖不足，导致部分含噪场景下泛化有限。第三，结果章已经补充本地200回合矩阵的距离、平滑度和置信区间指标，但这些指标仍来自仿真日志复算，不能直接代表真实机器人闭环性能。第四，完整闭环吞吐不等于真实控制实时性，策略推理、MPC求解和仿真开销需要拆分测试。

\section{未来工作}

后续工作可以从四个方向推进。第一，重新采集或增强噪声感知数据，使专家数据和学生策略暴露于更接近运行时的观测分布。第二，引入显式滤波或状态估计模块，缓解高斯抖动、漂移和延迟对参考构造的影响。第三，补充更完整的消融实验，包括目标中心化表示、未来条件、离散相位分支、动作块长度和时间集成温度等因素。第四，将策略推理延迟和MPC求解延迟分离评估，并在满足实时性后考虑向更高保真仿真或真实机器人平台迁移。

总体而言，本文验证了合成专家数据和扩散模仿策略在仿真动态抓取任务中的可行性，也揭示了学习策略与模型控制器在不同噪声类型下的互补失效模式。这为后续构建更鲁棒的动态目标机器人控制系统提供了实验基础。